% ============================================================================
\chapter{Asynchrone : boucle d'\'ev\'enements, promesses, \texttt{async/await}}
\label{ch:js-async}
% ============================================================================

\begin{objectif}
\begin{itemize}
  \item Expliquer ce qui se passe vraiment quand une requ\^ete est en cours.
  \item Cr\'eer, cha\^iner, combiner et annuler des promesses sans fuite ni \og ~promesse flottante~ \fg{}.
  \item \'Ecrire du code asynchrone lisible avec \verb|async/await|, en parall\`ele quand c'est possible.
  \item Impl\'ementer les motifs r\'eels~: retry, timeout, debounce, throttle, file \`a concurrence limit\'ee.
  \item Comprendre les it\'erateurs, g\'en\'erateurs et it\'erateurs asynchrones.
\end{itemize}
\end{objectif}

\section{Le mod\`ele d'ex\'ecution}

JavaScript est \textbf{mono-thread}~: un seul appel de fonction \`a la fois. Mais il est
\textbf{asynchrone et non bloquant}~: les op\'erations lentes (r\'eseau, disque, timer) sont d\'el\'egu\'ees
\`a des threads du syst\`eme (navigateur ou \link{libuv}{https://github.com/libuv/libuv} sous Node).
Leur r\'esultat est trait\'e quand la pile d'appel est vide.

\begin{enumerate}
  \item La \textbf{call stack} ex\'ecute le code synchrone jusqu'au bout.
  \item Les \textbf{API} (fetch, \verb|setTimeout|, fs) renvoient le contr\^ole imm\'ediatement ;
        la r\'eponse arrive plus tard, via un rappel plac\'e dans une file.
  \item La \textbf{microtask queue} (promesses r\'esolues, \verb|queueMicrotask|, MutationObserver)
        est vid\'ee \emph{enti\`erement} avant la suite~: chaque t\^ache a la priorit\'e.
  \item La \textbf{task queue} (timers, E/S, \'ev\'enements) attend que les microt\^aches soient vides,
        puis fournit \emph{une} t\^ache.
  \item \verb|requestAnimationFrame| s'ex\'ecute avant le rendu de l'image suivante.
\end{enumerate}

\begin{codejs}
console.log("1 synchrone");

setTimeout(() => console.log("6 timeout 0ms (macrotache)"), 0);
Promise.resolve().then(() => console.log("4 then"));
queueMicrotask(() => console.log("5 microtask"));
async function f() { console.log("2 debut async"); await null; console.log("3 apres await"); }
f();

// Ordre reel : 1, 2, 3, 4, 5, 6
// Regle mnemonique : tout le code synchrone, PUIS toutes les microtaches, PUIS une macrotache.
\end{codejs}

\begin{note}
\og ~JS est mono-thread donc il ne peut pas faire de calcul~ \fg{} est faux pour les calculs lourds~:
on les d\'eplace dans un \textbf{Web Worker} (navigueur) ou un \textbf{worker thread} (Node),
et on communique par messages. Voir \cref{sec:workers}.
\end{note}

\section{Rappels --- et pourquoi on les a quitt\'es}

\begin{codejs}
// Convention "error-first" heritee de Node : premier argument = erreur
function lireConfig(chemin, cb) {
  setTimeout(() => {
    if (!chemin) return cb(new Error("chemin requis"));
    cb(null, { chemin, ok: true });
  }, 30);
}
lireConfig("a.json", (err, res) => { if (err) throw err; /* ... */ });

// Le "pyramid of doom" est la vraie raison d'abandonner les callbacks imbriques
a(() => b(() => c(() => d(() => e(() => { /* on ne sait plus ou on est */ }))));
\end{codejs}

\section{Promesses}

Une promesse est un \textbf{objet qui repr\'esente une valeur pas encore disponible}, dans un des
trois \'etats~: \emph{pending}, \emph{fulfilled}, \emph{rejected} --- et qui ne change plus d'\'etat.

\begin{codejs}
const tache = new Promise((resolve, reject) => {
  setTimeout(() => resolve(42), 100);           // une seule resolution utile
});

tache
  .then((v) => v * 2)              // renvoyer une valeur -> nouvelle promesse resolue
  .then((v) => Promise.resolve(v + 1))   // renvoyer une promesse -> elle est "awaitee"
  .catch((e) => { console.error(e); throw new Error("echec traite"); })
  .finally(() => console.log("toujours"));   // ne transforme pas la valeur

// Les methodes de combinaison
const [u, posts] = await Promise.all([fetchUser(), fetchPosts()]);   // toutes ou echec
const resultats = await Promise.allSettled([p1, p2, p3]);  // toujours : [{status,value|reason}]
const premier = await Promise.race([requete, timeout]);            // le premier qui finit
const meilleur = await Promise.any([miroir1, miroir2]);            // le premier qui reussit
const p = Promise.withResolvers();            // {promise, resolve, reject} : propre a l'exterieur
const r = await Promise.try(() => fonctionPeutEtreSynchrone());    // n'importe quoi, uniformise

// await sur n'importe quoi (valeur, thenable, tableau)
const x = await 5;                 // 5
for (const v of await Promise.all(tab.map(async (i) => i * 2))) {}
\end{codejs}

\subsection{\'Ecrire des promesses correctes}

\begin{codejs}
// 1. Ne jamais creer une promesse "a la main" autour d'une API deja promise
const mauvais = new Promise((res) => res(fetch(url)));      // anti-pattern
const bon = fetch(url);

// 2. Promisify utile pour les API a callback
const promisifier = (fn) => (...args) =>
  new Promise((resolve, reject) =>
    fn(...args, (err, val) => (err ? reject(err) : resolve(val)))
  );
const tard = (ms) => new Promise((r) => setTimeout(r, ms));
await Promise.all([tard(100), tard(50)]);     // 100 ms (parallele), pas 150

// 3. Toujours gerter l'erreur, sinon : "unhandledRejection" (crash par defaut sous Node)
process.on("unhandledRejection", (reason) => { console.error("NON GREEE", reason); });

// 4. Une promesse n'est pas annulable : on ajoute un AbortController
const ctrl = new AbortController();
const reponse = await fetch(url, { signal: ctrl.signal });
ctrl.abort(new DOMException("Timeout", "TimeoutError"));  // -> fetch rejette

// 5. Enveloppe timeout reutilisable
async function avecTimeout(p, ms, message = "Timeout") {
  const ctrl = new AbortController();
  const t = setTimeout(() => ctrl.abort(new Error(message)), ms);
  try {
    return await Promise.race([
      p,
      new Promise((_, rej) => ctrl.signal.addEventListener("abort", () => rej(ctrl.signal.reason))),
    ]);
  } finally { clearTimeout(t); }
}
\end{codejs}

\section{\texttt{async/await}}

\verb|async function| renvoie toujours une promesse~; \verb|await| met la fonction en pause
et rend la main \`a la boucle d'\'ev\'enements.

\begin{codejs}
async function chargerProjet(id) {
  try {
    const rep = await fetch(`/api/projects/${id}`);
    if (!rep.ok) throw new Error(`HTTP ${rep.status} sur ${rep.url}`);
    const projet = await rep.json();
    const membres = await Promise.all(              // parallelisme : on attend les deux a la fois
      projet.membreIds.map((mid) => fetch(`/api/users/${mid}`).then((r) => r.json()))
    );
    return { ...projet, membres };
  } catch (err) {
    if (err.name === "AbortError") throw new Error("Requete annulee par l'utilisateur", { cause: err });
    throw err;                                       // remonter ce qu'on ne comprend pas
  } finally {
    masquerSpinner();
  }
}
\end{codejs}

\begin{tablez}[p{0.2\textwidth}XX]
  \headrow \thd{Anti-pattern} \& \thd{Probl\`eme} \& \thd{Correction} \\
  \midrule
  \verb|await p1; await p2;| (ind\'ependants) \& Temps cumul\'e, latence doubl\'ee \& \verb|await Promise.all([p1, p2])| \\
  \verb|tab.map(async x => await f(x))| \& concurrence non born\'ee (10 000 requ\^etes d'un coup) \& fichier d'attente born\'ee (ci-dessous) \\
  \verb|forEach(async x => ...)| \& la boucle ne \og ~attend~ \fg{} pas ; les erreurs sont perdues \& \verb|for...of| + \verb|await|, ou \verb|Promise.all(map)| \\
  \verb|.then(f).catch(g)| avec throw dans g \& \og ~unhandled rejection~ \fg{} \& attraper au plus pr\`es ou renvoyer \\
  \verb|catch {}| vide \& Bug invisible, donn\'ees corrompues \& au minimum \verb|console.error| ou \verb|throw| \\
  fonction non \verb|async| qui appelle une promesse \& valeur \verb|Promise \{pending\}| affich\'ee \& remonter \verb|async| jusqu'\`a l'appelant \\
\end{tablez}

\subsection{Trois utilitaires \`a avoir dans son projet}

\begin{codejs}
// 1. Retry avec recul exponentiel et jitter (la base de la resilience reseau)
async function retry(fn, { essais = 4, base = 200, max = 4000, surErreur } = {}) {
  let dernier;
  for (let i = 0; i < essais; i++) {
    try { return await fn(i); }
    catch (e) {
      dernier = e;
      surErreur?.(e, i);
      if (i === essais - 1) break;
      const attente = Math.min(max, base * 2 ** i) * (0.5 + Math.random() * 0.5);
      await new Promise((r) => setTimeout(r, attente));
    }
  }
  throw dernier;
}

// 2. Pool de concurrence : N taches en parallele, pas plus
async function parallelLimite(taches, limite = 5) {
  const resultats = new Array(taches.length);
  let curseur = 0;
  const worker = async () => {
    while (curseur < taches.length) {
      const i = curseur++;
      resultats[i] = await taches[i]();
    }
  };
  await Promise.all(Array.from({ length: Math.min(limite, taches.length) }, worker));
  return resultats;
}

// 3. Debounce / throttle : les deux fonctions les plus demandees en UI
const debounce = (fn, ms = 250) => {
  let t;
  const wrapped = (...args) => new Promise((res) => {
    clearTimeout(t);
    t = setTimeout(() => res(fn(...args)), ms);
  });
  wrapped.cancel = () => clearTimeout(t);
  return wrapped;
};
const throttle = (fn, ms = 100) => {
  let dernier = 0, minuterie = null, dernArgs;
  return (...args) => {
    const maintenant = Date.now();
    dernArgs = args;
    if (maintenant - dernier >= ms) { dernier = maintenant; fn(...args); }
    else if (!minuterie) minuterie = setTimeout(() => {
      dernier = Date.now(); minuterie = null; fn(...dernArgs);
    }, ms - (maintenant - dernier));
  };
};

const surRecherche = debounce(async (q) => { /* appel API */ }, 350);
document.querySelector("#q").addEventListener("input", (e) => surRecherche(e.target.value));
\end{codejs}

\section{G\'en\'erateurs, it\'erateurs et it\'erateurs asynchrones}

Un \textbf{it\'erateur} est un objet avec \verb|next()| renvoyant \verb|{ value, done }| ;
un \textbf{it\'erable} expose \verb|[Symbol.iterator]()|. \verb|for...of|, le spread et
\verb|Array.from| ne fonctionnent qu'avec \c{c}a.

\begin{codejs}
// Rendre un objet iterable (une playlist)
const playlist = {
  titres: ["a", "b", "c"],
  *[Symbol.iterator]() { for (const t of this.titres) yield t.toUpperCase(); },
};
[...playlist];                    // ["A","B","C"]

// Sequence infinie + paresse
function* entiersPairs(de = 0) { for (let n = de; ; n += 2) yield n; }
const dix = [];
for (const n of entiersPairs(2)) { if (dix.push(n) === 10) break; }

// Pipeline de transformation (generator-based)
function* map(iterable, fn) { for (const x of iterable) yield fn(x); }
function* filtre(iterable, pred) { for (const x of iterable) if (pred(x)) yield x; }
const lignes = [1, 2, 3, 4, 5, 6, 7, 8, 9, 10];
const sommeDesCarresImpairs = [...filtre(map(lignes, (x) => x * x), (x) => x % 2)]
  .reduce((a, b) => a + b, 0);

// Iterateur asynchrone : lire un tres gros fichier ligne a ligne sans le charger en RAM
async function* lignesFichier(url) {
  const res = await fetch(url);
  const decod = new TextDecoderStream("utf-8");
  const flux = res.body.pipeThrough(decod).pipeThrough(new TextEncoderStream ? new TransformStream() : new TransformStream());
  let tampon = "";
  for await (const chunk of res.body.pipeThrough(new TextDecoderStream("utf-8"))) {
    tampon += chunk;
    let idx;
    while ((idx = tampon.indexOf("\n")) >= 0) { yield tampon.slice(0, idx); tampon = tampon.slice(idx + 1); }
  }
  if (tampon) yield tampon;
}
// pour (const ligne of await ... ) n'existe pas : on ecrit
for await (const ligne of lignesFichier("https://exemple.gouv.fr/gros-fichier.csv")) { /* ... */ }
\end{codejs}

\subsection{Le flux (Streams) --- l'id\'ee \`a retenir}

Un \verb|ReadableStream| se lit morceau par morceau~; un \verb|WritableStream| s'\'ecrit de m\^eme ;
\verb|pipeThrough| et \verb|TransformStream| les cha\^inent. C'est ce qui permet de
t\'el\'echarger et transformer 5~Go sans jamais les mettre en m\'emoire, de streamer une r\'eponse
depuis Express ou un mod\`ele de langage, et de parser du SSE. On l'utilise concr\`etement dans
les chapitres Node et API.

\section{Chronologie, temps r\'eel et minuterie}

\begin{codejs}
const id = setInterval(() => tick(), 1000);   // recurrent
clearInterval(id);
let t = setTimeout(() => envoi(), 300);       // unique
clearTimeout(t);

// setTimeout(fn, 0) n'est pas "immediat" : c'est "prochaine macrotache"
// queueMicrotask(fn) est "des que la pile est vide" (avant le rendu)
// requestAnimationFrame(fn) est "juste avant le prochain affichage" (~16 ms)
function boucle() { dessiner(); requestAnimationFrame(boucle); }

// Mesurer proprement (pas Date.now)
const t0 = performance.now();
attente();
console.log(`${(performance.now() - t0).toFixed(2)} ms`);
\end{codejs}

\begin{exercice}
\begin{enumerate}
  \item \'Ecris \verb|attendre(ms)| puis transforme une fonction callback \verb|lireFichier| en
        version promise \verb|lireFichierAsync| (avec \verb|promisifier|~; puis \verb|fs/promises|).
  \item Pr\'edis l'ordre d'affichage, puis v\'erifie~:
\begin{codeplain}
console.log("A");
setTimeout(() => console.log("B"), 0);
Promise.resolve().then(() => console.log("C"));
console.log("D");
(async () => { console.log("E"); await null; console.log("F"); })();
console.log("G");
\end{codeplain}
  \item Charge les 26 lettres depuis \verb|https://api.npmjs.org/downloads/point/last-week/|
        combin\'ees avec une API locale~: utilise \verb|Promise.allSettled|, affiche le top 10 des
        paquets, et ne plante pas quand une requ\^ete \'echoue.
  \item Impl\'emente \verb|crawler(startUrl, maxPages)|~: il r\'ecup\`ere une liste pagin\'ee
        \verb|?page=n|, en parall\`ele \`a 5 maximum, s'arr\^ete sur page vide, ajoute un
        \verb|retry| \`a 3 essais et un \verb|timeout| de 5 s.
  \item \'Ecris un \verb|rateLimit(fn, appelsParSeconde)| qui r\'epartit les appels (token bucket),
        puis un \verb|pQueue(tasks, { concurrency, pause, resume })| minimal.
  \item Ajoute un bouton \og ~Annuler~ \fg{} \`a une recherche~: brancher un \verb|AbortController|,
        annuler la requ\^ete pr\'ec\'edente \`a chaque nouvelle frappe, et ignorer proprement
        \verb|AbortError|.
  \item \'Ecris une file d'attente de messages pour un chat~: \verb|emit|/\verb|on|, avec
        \verb|once|, \verb|off|, et un g\'en\'erateur asynchrone \verb|for await (const msg of bus)|.
\end{enumerate}
\end{exercice}

\begin{ressources}
\begin{itemize}
  \item MDN --- Promises --- \url{https://developer.mozilla.org/fr/docs/Web/JavaScript/Reference/Global\\\\_Objects/Promise}
  \item MDN --- Event loop et microtasks --- \url{https://developer.mozilla.org/en-US/docs/Web/API/HTML\\\\_DOM\\\\_API/Microtask\\\\_guide}
  \item Visualiseur interactif de la boucle d'\'ev\'enements --- \url{https://www.jsv9000.app/}
  \item Node --- Streams --- \url{https://nodejs.org/api/stream.html}
  \item Jake Archibald --- \emph{In The Loop} (la r\'ef\'erence sur l'ordre d'ex\'ecution) --- \url{https://jakearchibald.com/2015/events-based-runtimes/}
\end{itemize}
\end{ressources}

\index{JavaScript!asynchrone}\index{promesse}\index{async await}\index{event loop}
